\cotefile{34}
\scanpage{3}
\mspage{1}
\uline{Exposé I}\quad \deleted{\uline{Action du groupe de monodromie}} : \uline{Généralités}, \inserted{le} \uline{th. de monodromie par voie arithmétique}.

0. \uline{Introduction}

0.1. Notations \(S=\operatorname{Spec}\deleted{\unsure{R|A}}V\) un trait, \(s\) pt fermé, \(\eta\) point générique
\[
\begin{aligned}
\mathfrak m&=\pi V, & k&=k(s), & K&=k(\eta); & p&=\text{exp. car. de }k(s).\\[-2pt]
&\phantom{=}\shortparallel &&\phantom{=}\shortparallel\\[-2pt]
&\phantom{=}tV &&\phantom{=}V/\mathfrak m
\end{aligned}
\]
Le plus souvent, on suppose \(\deleted{\unsure{R|A}}V\) hensélien, et \(\overline{\mathrm m}\) strictement local i.e. \(k\) sép. clos ;

\deleted{\unsure{\(K=\)|Soit}} \(\overline K\) une clôture séparable de \(K\), \(\bar\eta=\operatorname{Spec}\overline K\) ;
\[
\pi=\operatorname{Gal}(\overline K/K).
\]

\(\ell\) nombre premier. \deleted{Dans la plus grande partie} \deleted{Sauf mention du contraire} \inserted{Pour tout ce \unsure{qu'on fera|qui se fera} de sérieux,} \deleted{\unsure{voir}} on suppose \deleted{L'objet \unsure{du|de}}
\[
\ell\ne\operatorname{car}k
\]

0.2. L'objet du Séminaire est l'étude de certaines représentations \(\ell\)-adiques de \(\pi\), provenant de diverses situations géométriques. Une repr. \(\ell\)-adique de \(\pi\) est un \deleted{\unsure{mor. de}} hom. continu\unsure{s}
\[
\pi\longrightarrow\operatorname{Aut}_{\mathbb Z_\ell}(E),\tag*{(*)}
\]
où \(E\) est un module de t.f. sur \(\mathbb Z_\ell\). On pourra aussi regarder
\[
\pi\longrightarrow\operatorname{Aut}_{\mathbb Q_\ell}(E)
\]
\(E\) module de t.f. sur \(\mathbb Q_\ell\).

\uline{Exemple} 0.3. \(X_K\) schéma algébrique sur \(K\), \inserted{\(i\) entier.} On suppose que \(H^i(X_{\overline K},\mathbb Z/\ell\mathbb Z)\) est un groupe fini, de sorte que
\[
E=H^i(X_{\overline K},\mathbb Z_\ell)\overset{\mathrm{dfn}}{=}\varprojlim_{\nu}H^i(X_{\overline K},\mathbb Z/\ell^{\nu}\mathbb Z)
\]
est un module de t.f. sur \(\mathbb Z_\ell\) ; \unsure{Cette|cette} hyp. \deleted{\illegible{1}} est \deleted{\unsure{K}} vérifiée dans les trois cas suivants :

a) On dispose de résolution des sing. \inserted{sous forme de Hironaka (  )} des schémas alg. de dim \(\leq\dim X\) sur \(\overline K\), -- p. ex. \(K\) de car. nulle

b) \(X\) propre sur \(K\).

c) \(X\) lisse sur \(K\).
% archive sheet number: 2
